\section*{Розділ I. Загальні положення}
\addcontentsline{toc}{section}{Розділ I. Загальні положення}

\subsection*{1.1. Статус Студентської ради студмістечка}
\addcontentsline{toc}{subsection}{1.1. Статус Студентської ради студмістечка}
    1.1.1. Студентська рада студмістечка (далі – СР СМ) є постійно діючим вищим колегіальним виконавчим органом студентського самоврядування Державного університету "Київський авіаційний інститут" (далі – Університет) у побутовій сфері.

    1.1.2. СР СМ діє відповідно до Конституції України, Закону України "Про вищу освіту", Статуту Університету, Положення про органи студентського самоврядування Університету (далі – Положення про ОСС), цього Положення та інших нормативно-правових актів України.

    1.1.3. Це Положення визначає повноваження, порядок формування, структуру, порядок роботи та інші аспекти діяльності СР СМ і діє відповідно до та на виконання Положення про ОСС.

\subsection*{1.2. Мета та основні завдання СР СМ}
\addcontentsline{toc}{subsection}{1.2. Мета та основні завдання СР СМ}
    1.2.1. Метою діяльності СР СМ є ефективне представництво та захист прав і законних інтересів студентів Університету у побутовій сфері, сприяння створенню належних умов проживання, відпочинку та соціально-побутового забезпечення у студентському містечку.

    1.2.2. Основними завданнями СР СМ є:

        \begin{enumerate}[label=\alph*)]
            \item Представлення інтересів студентів з побутових питань перед адміністрацією Університету та її структурними підрозділами, відповідальними за студентське містечко та гуртожитки.
            \item Забезпечення виконання рішень Конференції студентів Університету (КСУ) та реалізація завдань студентського самоврядування у побутовій сфері, визначених Положенням про ОСС.
            \item Координація діяльності студентських рад гуртожитків (СРГ).
            \item Сприяння створенню належних умов для проживання, відпочинку, безпеки та дозвілля студентів у гуртожитках та на території студмістечка.
            \item Участь у вирішенні питань щодо розподілу місць у гуртожитках, встановлення розміру плати за проживання, затвердження правил внутрішнього розпорядку гуртожитків та розвитку інфраструктури студмістечка.
            \item Сприяння соціальному забезпеченню та організації культурно-побутових заходів у студмістечку.
            \item Інформування студентів про діяльність ОСС у побутовій сфері та важливі аспекти життя студмістечка.
        \end{enumerate}

\subsection*{1.3. Принципи діяльності СР СМ}
\addcontentsline{toc}{subsection}{1.3. Принципи діяльності СР СМ}
    1.3.1. СР СМ здійснює свою діяльність на принципах, визначених Положенням про ОСС, зокрема: законності, добровільності, колегіальності, виборності, рівності прав студентів на участь у самоврядуванні, прозорості, організаційної самостійності.

    1.3.2. Ключовими принципами, що визначають виконавчий характер діяльності СР СМ, є підзвітність перед КСУ та ефективність у реалізації покладених завдань та повноважень.

\subsection*{1.4. Взаємозв'язок з іншими суб'єктами}
\addcontentsline{toc}{subsection}{1.4. Взаємозв'язок з іншими суб'єктами}
    1.4.1. СР СМ є підзвітною та підконтрольною Конференції студентів Університету.

    1.4.2. СР СМ взаємодіє з іншими органами студентського самоврядування (Студентською радою Київського авіаційного інституту (СР КАІ), Студентською уповноваженою делегацією (СУД), Центральною виборчою комісією студентів (ЦВКс), студентськими радами гуртожитків (СРГ)), адміністрацією Університету та студентськими організаціями у порядку, визначеному Положенням про ОСС та цим Положенням.

    1.4.3. СР СМ та СР КАІ є рівноправними паралельними виконавчими органами студентського самоврядування, кожен у своїй сфері компетенції, та взаємодіють на засадах партнерства і координації діяльності.
